\chapter{\(\mathrm{R}_{\mathrm{analytic}}\): Form, Difficulty, and the Riemann Hypothesis}
\label{v4-arith:R-analytic-comparison}

Chapter~\ref{v4-arith:hpar-det} isolates three inputs for a sufficient
argument toward the Riemann hypothesis: weight existence (W), the Weil
explicit formula at Gaussian test functions (H), and the sharp moment
asymptotic \(\mathrm{R}_{\mathrm{analytic}}\),
which asserts \(\mu_{2n}\asymp C^{n}n!^{2}\) unconditionally for the
symmetrised shifted-zero moment sequence of \(\xi\).
Remark~\ref{v4-arith:hpar-det:rem:residual-vs-RH} of
Chapter~\ref{v4-arith:hpar-det} characterises
\(\mathrm{R}_{\mathrm{analytic}}\) as weaker than RH in the implication
direction and as ``classical-analytic''. The statement has three
independent axes: form, implication
direction, and difficulty. The conclusions are (i) the form is
classical-analytic; (ii) this reduction does not derive RH from
\(\mathrm{R}_{\mathrm{analytic}}\) alone; (iii) the difficulty is not
classical but rather density-of-zeros, a major open problem adjacent to
RH in depth.
The Bombieri--Lagarias upper bound does not imply Carleman
divergence. After the \(1/(2n)\)-th root and inversion, it yields the
convergent Bertrand comparison \(\sum 1/(n(\log n)^2)\). This is why
Proposition~\ref{v4-arith:hpar-det:prop:BL-carleman-partial}
records the Bombieri--Lagarias construction as inconclusive for the
determinacy estimate rather than as a divergence proof.

\section{The Four Assertions}
\label{v4-arith:r-comparison:statement}

\begin{proposition}[form of \(\mathrm{R}_{\mathrm{analytic}}\)]
\label{v4-arith:r-comparison:claim-1}
\(\mathrm{R}_{\mathrm{analytic}}\) is classical-analytic in form:
a two-sided moment asymptotic
\(C_{-}^{n}n!^{2}\le \mu_{2n}\le C_{+}^{n}n!^{2}\) on a regularised
positive measure associated to the symmetrised shifted zeros of
\(\xi\).
\end{proposition}

\begin{proposition}[no direct implication to RH]
\label{v4-arith:r-comparison:claim-2}
\(\mathrm{R}_{\mathrm{analytic}}\) is weaker than RH in the conclusion it
states: an off-line-zero configuration with subleading \(2n\)-th moment
contribution is not excluded by the moment statement alone. Thus this
reduction does not prove
\(\mathrm{R}_{\mathrm{analytic}}\Rightarrow\mathrm{RH}\).
(\Cref{v4-arith:hpar-det:rem:residual-vs-RH}.)
\end{proposition}

\begin{proposition}[Bombieri--Lagarias gives a finite Bertrand comparison]
\label{v4-arith:r-comparison:claim-3}
\ClaimStatusProvedHere. The upper bound
\[
\mu_{2n}\le (2\pi)^{-4n}(2n)!(\log n)^{4n}(1+o(1))
\]
implies only
\[
\mu_{2n}^{-1/(2n)}\geq
\frac{c}{n(\log n)^2}(1+o(1))
\]
for some \(c>0\). The comparison series
\(\sum_{n\geq 2}1/(n(\log n)^2)\) converges. No Carleman divergence
follows from this estimate.
\end{proposition}

\begin{proposition}[form, implication, difficulty]
\label{v4-arith:r-comparison:claim-paraphrase}
\(\mathrm{R}_{\mathrm{analytic}}\) has classical-analytic form and is
not an RH proof by itself. Its difficulty is nevertheless RH-adjacent:
removing the logarithmic inflation in the moment asymptotic requires
density-of-zeros input beyond the Bombieri--Lagarias and Conrey
estimates used here.
\end{proposition}

\section{The Stirling Scale in the Carleman Computation}
\label{v4-arith:r-comparison:stirling-error}

\Cref{v4-arith:r-comparison:claim-3} follows from the
Stirling asymptotic \((2n)!^{1/(2n)}\sim 2n/e\). The power grows
linearly in \(n\). The factor \((\log n)^{4n}\) in the upper bound
on \(\mu_{2n}\) becomes \((\log n)^{+2}\) after taking the
\(1/(2n)\)-th power of \(\mu_{2n}\), hence \((\log n)^{-2}\) in
\(\mu_{2n}^{-1/(2n)}\). The sign of the logarithmic exponent reverses
under inversion, giving the convergent Bertrand series
\(\sum 1/(n(\log n)^{2})\).

\section{Bombieri--Lagarias Bounds \(\lambda_{n}\), Not \(\mu_{2n}\)}
\label{v4-arith:r-comparison:BL-lambda-not-mu}

Bombieri--Lagarias~1999
(\Cref{v4-arith:vb:thm:BL-bound}) bounds the Li coefficient
\(\lambda_{n}=\sum_{\rho}[1-(1-1/\rho)^{n}]\) by
\(\lambda_{n}\leq n(c_{0}+\log n)/2\). This is a bound on a highly
oscillatory unit-modulus phase sum: under RH, \(|1-1/\rho|=1\) for
every non-trivial zero \(\rho\), so \(\lambda_{n}\) grows
sub-polynomially via phase cancellation. The moment sum
\(\mu_{2n}=\sum_{\rho}(1/4+\gamma_{\rho}^{2})^{n}\) is a monotone
positive sum whose growth is governed by the density of zeros high
in the critical strip via Riemann--von Mangoldt, with no
cancellation. These two quantities are related only through
\(\mathrm{H\text{-}P}^{\mathrm{Ar}}\)-conditional spectral-trace
identities (\Cref{v4-arith:hpar-det:cycle-5}), not through
elementary norm bounds; a direct transfer of the Bombieri--Lagarias
bound to \(\mu_{2n}\) requires a phase-cancellation-to-magnitude
comparison that is not available. The upper bound on
\(\mu_{2n}\) in \Cref{v4-arith:hpar-det:prop:BL-carleman-partial}
comes from Riemann--von Mangoldt zero counting combined with
RH-conditional saddle-point evaluation, not from the
Bombieri--Lagarias Li-coefficient bound. Closing the
\((\log n)^{4n}\to 1\) gap requires density-of-zeros improvements
on the critical line beyond Conrey's 40\%, via the
Selberg--Levinson--Conrey--Feng density theory.

\section{Position of \(\mathrm{R}_{\mathrm{analytic}}\) in the Equivalence Chain}
\label{v4-arith:r-comparison:chain-position}

Chapter~\ref{v4-arith:HPAr-vbconverse}
(\Cref{v4-arith:hpvbconv:thm:VBstar-converse-spectral-regularity})
establishes the four-way equivalence
\[
(\mathrm{VB^{*}c})\;\Leftrightarrow\;(\mathrm{Det})\;\Leftrightarrow\;
(\mathrm{HB+})\;\Leftrightarrow\;(\mathrm{BD+})
\]
under P2, P3, and \(\mathrm{H\text{-}P}^{\mathrm{Ar}}\), where
\((\mathrm{VB^{*}c})\) is the \emph{converse} ``zero-placement
\(\Rightarrow\) \(\mathrm{VB}^{*}\)'', not the main direction
\(\mathrm{VB}^{*}\Rightarrow\) zero-placement. Determinacy is equivalent to
the VB\(^{*}\) \emph{converse}. Chapter~\ref{v4-arith:HPAr-vbconverse}
further records
(\Cref{v4-arith:hpvbconv:rem:strict-content},
\Cref{v4-arith:hpvbconv:cor:three-strictly-finer}) that all four
conditions require input beyond RH in this reduction: in particular,
no implication \(\mathrm{RH}\Rightarrow\mathrm{Det}\) is proved here,
and \(\mathrm{Det}\Rightarrow\mathrm{RH}\) requires
\(\mathrm{H\text{-}P}^{\mathrm{Ar}}\) plus the VB\(^{*}\) main direction
(\Cref{v4-arith:vb:thm:three-equivalences}). Under P2, P3,
\(\mathrm{H\text{-}P}^{\mathrm{Ar}}\), one has
\(\mathrm{R}_{\mathrm{analytic}}\Rightarrow\mathrm{Det}
\Rightarrow(\mathrm{VB^{*}c})\), placing
\(\mathrm{R}_{\mathrm{analytic}}\) as an input to the converse
direction, not to the main direction and not to RH directly.
Chapter~\ref{v4-arith:hpar-det}'s characterisation of
\(\mathrm{R}_{\mathrm{analytic}}\) as not RH-equivalent in the
implication direction is therefore correct, but the reason is not
that the statement is classical: it is that determinacy requires
additional input in the implication hierarchy, as established by
Chapter~\ref{v4-arith:HPAr-vbconverse}.

\section{Carleman Sum from the Bombieri--Lagarias Upper Bound}
\label{v4-arith:r-comparison:carleman-BL-explicit}

\subsection{Upper Bound on \(\mu_{2n}\) via Riemann--von Mangoldt}
\label{v4-arith:r-comparison:upper-bound}

Proposition~\ref{v4-arith:hpar-det:prop:BL-carleman-partial} proves,
via Riemann--von Mangoldt zero-counting and saddle-point evaluation,
\begin{equation}
\label{v4-arith:r-comparison:eq:upper-bound}
\mu_{2n}\;\leq\;\frac{(2n)!}{(2\pi)^{4n}}\cdot(\log n)^{4n}\cdot(1+o(1)).
\end{equation}
The saddle-point derivation is standard and is not disputed here.

\subsection{Stirling Evaluation}
\label{v4-arith:r-comparison:corrected-stirling}

\begin{lemma}[Stirling for \((2n)!^{1/(2n)}\)]
\label{v4-arith:r-comparison:lem:stirling}
\ClaimStatusProvedHere. By Stirling's formula
\(n!=\sqrt{2\pi n}\,(n/e)^{n}(1+O(1/n))\),
\begin{equation}
(2n)!^{1/(2n)}\;=\;\frac{2n}{e}\cdot(4\pi n)^{1/(4n)}\cdot(1+O(1/n)).
\end{equation}
Since \((4\pi n)^{1/(4n)}\to 1\) as \(n\to\infty\), one has
\((2n)!^{1/(2n)}\sim 2n/e\), i.e., the \((2n)!^{1/(2n)}\) power
grows \emph{linearly} in \(n\), not as \(\sqrt{n}\).
\end{lemma}

\begin{proof}
\((2n)!=\sqrt{4\pi n}\,(2n/e)^{2n}(1+O(1/n))\) gives
\begin{align*}
(2n)!^{1/(2n)}
&=\bigl[\sqrt{4\pi n}\,(2n/e)^{2n}\bigr]^{1/(2n)}\cdot(1+O(1/n))^{1/(2n)}\\
&=(4\pi n)^{1/(4n)}\cdot(2n/e)\cdot\bigl(1+O(1/n^{2})\bigr).
\end{align*}
Since \((4\pi n)^{1/(4n)}=\exp(\tfrac{1}{4n}\log(4\pi n))
=1+\tfrac{\log(4\pi n)}{4n}+O(n^{-2}(\log n)^{2})\to 1\), the claim
follows.
\end{proof}

\begin{proposition}[Carleman-sum lower bound from the upper bound]
\label{v4-arith:r-comparison:prop:BL-carleman-corrected}
\ClaimStatusProvedHere. From
\Cref{v4-arith:r-comparison:eq:upper-bound},
\begin{equation}
\label{v4-arith:r-comparison:eq:mu-power}
\mu_{2n}^{1/(2n)}\;\leq\;
\frac{1}{(2\pi)^{2}}\cdot(2n)!^{1/(2n)}\cdot(\log n)^{2}\cdot(1+o(1))
\;\sim\;\frac{2n}{e(2\pi)^{2}}\cdot(\log n)^{2}.
\end{equation}
Hence
\begin{equation}
\label{v4-arith:r-comparison:eq:mu-inverse-lower}
\mu_{2n}^{-1/(2n)}\;\geq\;
\frac{e(2\pi)^{2}}{2n(\log n)^{2}}\cdot(1-o(1))
\;=\;\frac{c_{\mathrm{BL}}}{n(\log n)^{2}}
\end{equation}
for a constant \(c_{\mathrm{BL}}>0\). Summing,
\begin{equation}
\label{v4-arith:r-comparison:eq:carleman-bound}
\sum_{n\geq 2}\mu_{2n}^{-1/(2n)}\;\geq\;
c_{\mathrm{BL}}\cdot\sum_{n\geq 2}\frac{1}{n(\log n)^{2}}.
\end{equation}
\end{proposition}

\begin{proof}
Direct substitution of \Cref{v4-arith:r-comparison:lem:stirling} into
the \(1/(2n)\)-th power of the upper bound
\Cref{v4-arith:r-comparison:eq:upper-bound}. The factor \((\log n)^{4n}\)
inside the bound becomes \((\log n)^{+2}\) after the
\(1/(2n)\)-th-root, and when inverted for \(\mu_{2n}^{-1/(2n)}\) becomes
\((\log n)^{-2}\). The factor \((2\pi)^{-4n}\) becomes \((2\pi)^{-2}\)
and when inverted becomes \((2\pi)^{+2}\). The factor
\((2n)!^{1/(2n)}\sim 2n/e\) in the numerator of \(\mu_{2n}^{1/(2n)}\)
becomes \(e/(2n)\) in \(\mu_{2n}^{-1/(2n)}\).
\end{proof}

\begin{corollary}[the BL Carleman construction is \emph{inconclusive}, not divergent]
\label{v4-arith:r-comparison:cor:inconclusive}
\ClaimStatusProvedHere. The lower bound on the Carleman sum from
\Cref{v4-arith:r-comparison:prop:BL-carleman-corrected} is
\begin{equation}
\label{v4-arith:r-comparison:eq:bertrand}
\sum_{n\geq 2}\frac{1}{n(\log n)^{2}}\;<\;+\infty,
\end{equation}
a classical Bertrand series which \emph{converges}. Hence the
lower bound on \(\sum_{n}\mu_{2n}^{-1/(2n)}\) produced by the BL
upper bound on \(\mu_{2n}\) is finite. This does \emph{not} prove
the Carleman sum diverges; it is strictly inconclusive.
\end{corollary}

\begin{proof}
The Bertrand series \(\sum_{n\geq 2}1/(n(\log n)^{\alpha})\)
converges iff \(\alpha>1\), by the integral comparison
\(\int_{2}^{\infty}dx/(x(\log x)^{\alpha})
=(\alpha-1)^{-1}(\log 2)^{-(\alpha-1)}\) for \(\alpha>1\). For
\(\alpha=2\) this integral is finite, confirming convergence.

A lower bound on \(\sum \mu_{2n}^{-1/(2n)}\) which is itself finite
does not establish divergence of \(\sum \mu_{2n}^{-1/(2n)}\). The
Carleman sum may still diverge or converge; the BL upper bound on
\(\mu_{2n}\) with the \((\log n)^{4n}\) log-factor is simply too
weak to decide.
\end{proof}

\begin{remark}[Stirling scale in \Cref{v4-arith:r-comparison:claim-3}]
\label{v4-arith:r-comparison:rem:delta-error}
\Cref{v4-arith:r-comparison:claim-3} asserts
\(\sum \mu_{2n}^{-1/(2n)}\geq c'\sum 2\pi(\log n)/\sqrt{2n}\). This
expression has \(\sqrt{n}\) in the denominator instead of \(n\), and
\((\log n)^{+1}\) in the numerator instead of \((\log n)^{-2}\). The
resulting sum \(\sum (\log n)/\sqrt{n}\) diverges, but the
transcription from \(\mu_{2n}^{-1/(2n)}\) via Stirling should
produce \(1/(n(\log n)^{2})\) per \Cref{v4-arith:r-comparison:lem:stirling}
and power-rearrangement of \((\log n)^{4n}\), giving a convergent
Bertrand series. The source is the approximation
\((2n)!^{1/(2n)}\sim\sqrt{2n}\) instead of \(\sim 2n/e\). Stirling's
formula gives \(n!^{1/n}\sim n/e\).
\end{remark}

\begin{remark}[Carleman divergence under the RH-heuristic asymptotic \(\mu_{2n}\asymp C^{n}(n!)^{2}\)]
\label{v4-arith:r-comparison:rem:heuristic-saves}
If one substitutes the RH-heuristic two-sided asymptotic
\(\mu_{2n}\asymp C^{n}(n!)^{2}\) with no \((\log n)^{4n}\) factor, as in
\Cref{v4-arith:hpvbconv:prop:carleman-insufficient}, then
\(\mu_{2n}^{1/(2n)}\asymp C^{1/2}\cdot n/e\) giving
\(\mu_{2n}^{-1/(2n)}\asymp e/(C^{1/2}n)\) and the Carleman sum
\emph{does} diverge at the harmonic rate \(\sum 1/n\). But this
heuristic is \emph{not} established unconditionally; the
BL-plus-Conrey machinery has the \((\log n)^{4n}\) gap
(\Cref{v4-arith:hpar-det:cor:unconditional-gap}) which puts the
true sum into the log-convergent regime of
\Cref{v4-arith:r-comparison:cor:inconclusive}. Hence Carleman's
criterion from BL-plus-Conrey alone is inconclusive; the
classical-heuristic form is what requires unconditional
confirmation.
\end{remark}

\section{Sharpening the \((\log n)^{4n}\) Gap Requires Zero-Density Hypotheses}
\label{v4-arith:r-comparison:gap-analysis}

\subsection{The Bombieri--Lagarias Polylog Bound Does Not Sharpen to \(\mu_{2n}\asymp n!^{2}\) by Classical Analysis}
\label{v4-arith:r-comparison:can-BL-sharpen}

\begin{proposition}[Bombieri--Lagarias bounds \(\lambda_{n}\); \(\mu_{2n}\) is governed by Riemann--von Mangoldt]
\label{v4-arith:r-comparison:prop:BL-vs-mu}
\ClaimStatusProvedHere. Bombieri--Lagarias Theorem~2
(\Cref{v4-arith:vb:thm:BL-bound}) bounds
\(\lambda_{n}=\sum_{\rho}[1-(1-1/\rho)^{n}]\), a heavily
oscillatory unit-modulus-phase sum in \(\rho\). Under RH, \(|1-1/\rho|=1\)
for every non-trivial zero, so the Bombieri--Lagarias bound is a
statement about phase-cancellation. The moment sum
\(\mu_{2n}=\sum_{\rho}(1/4+\gamma_{\rho}^{2})^{n}\) is a monotone
positive sum (under RH) whose growth is governed by the density of
zeros high in the critical line via Riemann--von Mangoldt. The two
quantities are related through
\(\mathrm{H\text{-}P}^{\mathrm{Ar}}\)-conditional spectral-trace
identities, not through elementary norm bounds.
\end{proposition}

\begin{proof}
The Bombieri--Lagarias bound gives
\(\lambda_{n}\geq-(n/2)\log n-O(n)\)
(\Cref{v4-arith:vb:thm:BL-bound}): polylog growth. A polylog bound
on \(\lambda_{n}\) carries no information about \(\mu_{2n}\):
\(\lambda_{n}\) grows via phase cancellation while \(\mu_{2n}\) is a
monotone positive sum with no cancellation. The identity
\(\mathrm{tr}(T_{n})=n\lambda_{n}\) holds only at the trace level
under \(\mathrm{H\text{-}P}^{\mathrm{Ar}}\)
(\Cref{v4-arith:hpar-det:cycle-5}).
\end{proof}

\begin{proposition}[sharpening the unconditional \(\mu_{2n}\) upper bound from
\((\log n)^{4n}\) to \((\log n)^{o(n)}\) requires zero-density improvements]
\label{v4-arith:r-comparison:prop:closing-gap-is-density-RH}
\ClaimStatusProvedHere. The \((\log n)^{4n}\) factor in
\Cref{v4-arith:r-comparison:eq:upper-bound} comes from the
Conrey-style fluctuation envelope in the RvM zero-counting:
\(N(T)\sim (T/2\pi)\log(T/2\pi)\) with second-order error
\(O(\log T)\), so that the integral representation of \(\mu_{2n}\)
picks up a \((\log T)^{\cdot}\) factor at the saddle cutoff
\(T\sim\sqrt{n}\log n\). Removing this factor to obtain
\(\mu_{2n}\leq C^{n}n!^{2}\) \emph{with no logarithmic inflation}
requires improving the zero-counting asymptotic from the
unconditional \(N(T)=(T/2\pi)\log(T/2\pi)-T/2\pi+O(\log T)\) to a
form where the error term does not amplify by a factor of \(\log T\)
at each moment-integral saddle. This is a zero-density-hypothesis
improvement, a classical analytic-number-theory problem that is
currently open and is not known to follow from RH alone.
\end{proposition}

\begin{proof}
The RvM estimate \(N(T)=(T/2\pi)\log(T/2\pi)-T/2\pi+O(\log T)\)
(Edwards~1974 Thm.~6.7) has error \(O(\log T)\). In the saddle-point
evaluation \(\mu_{2n}\approx \int_{0}^{\infty}t^{2n}(dN/dt)dt\), the
saddle is at \(t\sim\sqrt{n}\log n\). The error term \(O(\log T)\)
evaluated at the saddle contributes a factor of
\(O(\log(\sqrt{n}\log n))\sim O(\log n)\) at the saddle, and the
\(n\)-fold application (from the \(n\)-th moment) amplifies this to
\(O((\log n)^{4n})\) in the final bound (the factor of 4 comes from
the \(\mu_{2n}\) being the \(2n\)-th moment, with two such
amplifications in the saddle-point evaluation and two from the
\(t^{2}\) factor).

The relevant fluctuation is
\(S(T)=\tfrac{1}{\pi}\arg\zeta(\tfrac12+iT)\), the oscillating part of
\(N(T)=\tfrac{T}{2\pi}\log\tfrac{T}{2\pi e}+\tfrac78+S(T)+O(1/T)\);
unconditionally \(S(T)=O(\log T)\) (von Mangoldt--Backlund), and this is
precisely the error that the \(n\)-fold saddle amplifies. Reducing the
amplified contribution to \((\log n)^{o(n)}\) cannot be achieved by
sharpening this pointwise bound: \(S(T)\) is provably unbounded
(Selberg~1946), so no \(O(1)\) error term for \(N(T)\) exists. Removing
the inflation instead requires zero-density control on the zeros high
in the critical strip (a density-hypothesis-type input in the sense
of Bombieri~1965) or a cancellation in the regularised moment
integral that the crude magnitude bound discards. Either is a major
open analytic-number-theory problem, adjacent to RH in depth.
\end{proof}

\begin{remark}[classical-analytic form, RH-adjacent difficulty]
\label{v4-arith:r-comparison:rem:form-vs-difficulty}
\Cref{v4-arith:r-comparison:prop:closing-gap-is-density-RH} establishes
the key distinction: the \emph{form} of
\(\mathrm{R}_{\mathrm{analytic}}\) is classical-analytic (a moment
asymptotic), but its \emph{difficulty} is RH-adjacent. The gap
between BL-Conrey's polylog-inflated bound and the target
\(C^{n}n!^{2}\) bound is precisely the density-of-zeros hypothesis for
\(\zeta\), which is a major open problem on the same level as RH.
Density-of-zeros is weaker in conclusion than RH: a density-one
statement controls the exceptional set but does not put every zero on
the line. Thus \(\mathrm{R}_{\mathrm{analytic}}\) is not an RH proof in
this reduction, and it is not known unconditionally either.
\end{remark}

\subsection{Three Independent Axes: Form, Implication, Difficulty}
\label{v4-arith:r-comparison:axis-distinction}

\Cref{v4-arith:r-comparison:claim-paraphrase}'s conflation of classical
form with classical difficulty is resolved by separating three axes.

\begin{description}
\item[Form.] The statement of \(\mathrm{R}_{\mathrm{analytic}}\) is
the moment asymptotic
\(C_{-}^{n}n!^{2}\leq\mu_{2n}\leq C_{+}^{n}n!^{2}\): an object
in classical analysis, a moment sequence controlled by
saddle-point asymptotics.
\item[Implication direction.] The present reduction does not prove
\(\mathrm{R}_{\mathrm{analytic}}\Rightarrow\mathrm{RH}\): off-line zeros
with subleading \(2n\)-th moment contribution are not excluded by the
moment statement alone. The converse
\(\mathrm{RH}\Rightarrow\mathrm{R}_{\mathrm{analytic}}\) holds
partially under regularisation: Riemann--von Mangoldt plus
saddle-point gives \(\mu_{2n}\asymp (2\pi)^{-4n}(2n)!\) up to
logarithmic corrections, which by Stirling is \(C^{n}n!^{2}\) up
to the density-hypothesis refinements required to remove the
\((\log n)^{4n}\) factor.
\item[Difficulty.] The remaining gap
(\Cref{v4-arith:r-comparison:prop:closing-gap-is-density-RH}) is the
density-of-zeros hypothesis for \(\zeta\), a major open problem
of analytic number theory at the level of the density-1 hypothesis,
weaker than RH in its conclusion but adjacent to RH in depth.
\end{description}

\begin{center}
\begin{tabular}{|l|l|l|l|}
\hline
 & \textbf{Form} & \textbf{Implication to RH} & \textbf{Difficulty} \\
\hline
\(\mathrm{R}_{\mathrm{analytic}}\) & Classical-analytic & Does not imply RH & RH-adjacent \\
 & (moment asymptotic) & (weaker than RH in & (density-of-zeros \\
 & & one direction) & open problem) \\
\hline
\end{tabular}
\end{center}

\Cref{v4-arith:r-comparison:claim-paraphrase} infers from the classical
form that the difficulty is classical; this conflates the first and
third axes. The implication axis (\Cref{v4-arith:r-comparison:claim-2})
is correct: \(\mathrm{R}_{\mathrm{analytic}}\) is strictly weaker
than RH in the implication direction. The difficulty axis is
RH-adjacent.

\section{Determinacy, \(\mathrm{VB}^{*}\), and RH}
\label{v4-arith:r-comparison:chain-comparison}

The two propositions fix the two equivalences in the chain
\(\mathrm{Det}\Leftrightarrow\mathrm{VB}^{*}\text{ main direction}\Leftrightarrow\mathrm{RH}\).

\subsection{Step 1: Determinacy \(\Leftrightarrow\) VB\(^{*}\) (Which Direction?)}
\label{v4-arith:r-comparison:step-1}

\begin{proposition}[Determinacy is equivalent to the \(\mathrm{VB}^{*}\) converse]
\label{v4-arith:r-comparison:prop:det-converse-not-main}
\ClaimStatusProvedHere.
\Cref{v4-arith:hpvbconv:thm:VBstar-converse-spectral-regularity}
asserts
\[
(\mathrm{VB^{*}c})\;\Leftrightarrow\;(\mathrm{Det})\;\Leftrightarrow\;
(\mathrm{HB+})\;\Leftrightarrow\;(\mathrm{BD+})
\]
under P2, P3, \(\mathrm{H\text{-}P}^{\mathrm{Ar}}\), where
\((\mathrm{VB^{*}c})\) is ``zero-placement \(\Rightarrow\) \(\mathrm{VB}^{*}\)''.
The main direction \(\mathrm{VB}^{*}\Rightarrow\) zero-placement is the content
of \Cref{v4-arith:vb:thm:three-equivalences}, a separate conditional
implication not equivalent to determinacy.
\end{proposition}

\begin{proof}
\Cref{v4-arith:hpvbconv:rem:converse-does-not-help} records:
``The forward direction \(\mathrm{VB}^{*}\Rightarrow\mathrm{F\text{-}RH}\)
is the separate content of \Cref{v4-arith:rh:thm:VB-implies-F-RH},
which is a conditional theorem under P2, P3,
\(\mathrm{H\text{-}P}^{\mathrm{Ar}}\), and a positive Bost--Connes
regularised trace. The four-way equivalence of the converse does not
close the forward implication.'' Determinacy is equivalent to the
converse direction.
\end{proof}

\subsection{Step 2: Determinacy vs RH}
\label{v4-arith:r-comparison:step-2}

\begin{proposition}[Determinacy needs input beyond RH in this reduction]
\label{v4-arith:r-comparison:prop:det-strictly-finer}
\ClaimStatusProvedHere.
\Cref{v4-arith:hpvbconv:cor:three-strictly-finer} asserts that each
of \(\mathrm{VB}^{*}\)-converse, determinacy, and
Ba\'ez-Duarte sharp decay carries input beyond RH in this reduction.
The implication \(\mathrm{Det}\Rightarrow\mathrm{RH}\) is used only
with \(\mathrm{H\text{-}P}^{\mathrm{Ar}}\)
(\Cref{v4-arith:hpvbconv:thm:HPAr-implies-reduction}), and no
implication \(\mathrm{RH}\Rightarrow\mathrm{Det}\) is proved here.
\end{proposition}

\begin{proof}
\Cref{v4-arith:hpvbconv:rem:strict-content}: ``No structural
property of the zero distribution of \(\xi\) independent of RH is
known to force determinacy. Hence the \(\mathrm{VB}^{*}\) converse
requires an additional arithmetic input beyond RH.'' Hence this
chapter does not derive \(\mathrm{Det}\) from RH. The reverse implication
\(\mathrm{Det}\Rightarrow\mathrm{RH}\) requires
\(\mathrm{H\text{-}P}^{\mathrm{Ar}}\) plus the \(\mathrm{VB}^{*}\)
main direction (\Cref{v4-arith:vb:thm:three-equivalences}); the
two inputs are kept separate without auxiliary hypotheses.
\end{proof}

\subsection{The Implication Diagram}
\label{v4-arith:r-comparison:corrected-chain}

\begin{remark}[Position of \(\mathrm{R}_{\mathrm{analytic}}\) in the implication hierarchy]
\label{v4-arith:r-comparison:diag:chain}
Under P2, P3, \(\mathrm{H\text{-}P}^{\mathrm{Ar}}\):
\begin{center}
\(\mathrm{R}_{\mathrm{analytic}}\;\Longrightarrow\;\text{(Carleman criterion)}\;\Longrightarrow\;\mathrm{Det}\;\Longleftrightarrow\;(\mathrm{VB}^{*}\text{ converse})\)
\end{center}
Independently, as a conditional theorem
(\Cref{v4-arith:vb:thm:three-equivalences}):
\begin{center}
\(\mathrm{VB}^{*}\text{ main direction}\;\Longrightarrow\;\mathrm{RH}.\)
\end{center}
Combined (\Cref{v4-arith:hpvbconv:thm:HPAr-implies-reduction}):
\begin{center}
\(\mathrm{H\text{-}P}^{\mathrm{Ar}}+\mathrm{R}_{\mathrm{analytic}}+\mathrm{VB}^{*}\text{ main direction}\;\Longrightarrow\;\mathrm{VB}^{*}\;\Longleftrightarrow\;\mathrm{RH}+\mathrm{Det}.\)
\end{center}
\end{remark}

\begin{remark}[\(\mathrm{R}_{\mathrm{analytic}}\) controls the converse]
\label{v4-arith:r-comparison:rem:closes-converse}
In the implication diagram, \(\mathrm{R}_{\mathrm{analytic}}\) contributes
to the VB\(^{*}\) \emph{converse} direction (RH\(\Rightarrow\)VB\(^{*}\)),
by forcing Hamburger determinacy which together with H-P\(^{\mathrm{Ar}}\)
and RH gives VB\(^{*}\). \(\mathrm{R}_{\mathrm{analytic}}\) does not
by itself imply the VB\(^{*}\) \emph{main} direction nor RH. The
main direction requires VB\(^{*}\) proper or
some alternative construction to the Li positivity, which is not supplied
by \(\mathrm{R}_{\mathrm{analytic}}\).
\end{remark}

\section{Conclusion}
\label{v4-arith:r-comparison:conclusion}

\begin{theorem}[four assertions on \(\mathrm{R}_{\mathrm{analytic}}\)]
\label{v4-arith:r-comparison:thm:conclusion}
\ClaimStatusProvedHere. The four assertions
\Cref{v4-arith:r-comparison:claim-1}--\Cref{v4-arith:r-comparison:claim-paraphrase}
have the following mathematical form.
\begin{enumerate}
\item[\textbf{(1)}] \emph{Form}
(\Cref{v4-arith:r-comparison:claim-1}): \(\mathrm{R}_{\mathrm{analytic}}\)
has classical-analytic form, a two-sided moment asymptotic.
\item[\textbf{(2)}] \emph{Implication}
(\Cref{v4-arith:r-comparison:claim-2}):
\(\mathrm{R}_{\mathrm{analytic}}\) is not an RH proof by itself in this
reduction. The converse
\(\mathrm{RH}\Rightarrow\mathrm{R}_{\mathrm{analytic}}\) also
requires density-of-zeros refinements beyond RH alone.
\item[\textbf{(3)}] \emph{Bombieri--Lagarias scale}
(\Cref{v4-arith:r-comparison:claim-3}): the Bombieri--Lagarias--Conrey
construction gives a Carleman sum lower bound
\(\sum 1/(n(\log n)^{2})\), a Bertrand series
that converges, hence the Bombieri--Lagarias construction is inconclusive,
not divergent.
\item[\textbf{(4)}] \emph{Difficulty}
(\Cref{v4-arith:r-comparison:claim-paraphrase}): ``classical-analytic,
not RH-equivalent'' is correct on the form and implication axes. Its
difficulty is RH-adjacent, because the logarithmic inflation can only
be removed here by density-of-zeros input.
\end{enumerate}
\end{theorem}

\begin{proof}
(1) is \Cref{v4-arith:r-comparison:axis-distinction}: the statement is a
moment asymptotic.

(2) is \Cref{v4-arith:hpar-det:rem:residual-vs-RH} combined with
the off-line-zeros robustness argument.

(3) is \Cref{v4-arith:r-comparison:prop:BL-carleman-corrected} and
\Cref{v4-arith:r-comparison:cor:inconclusive}: the Stirling identity
\((2n)!^{1/(2n)}\sim 2n/e\) gives the convergent Bertrand series
\(\sum 1/(n(\log n)^{2})\).

(4) is \Cref{v4-arith:r-comparison:rem:form-vs-difficulty} and
\Cref{v4-arith:r-comparison:prop:closing-gap-is-density-RH}: removing the
logarithmic gap requires density-of-zeros improvements which,
while weaker than RH in conclusion, constitute a major open problem of
analytic number theory.
\end{proof}

\begin{corollary}[\(\mathrm{R}_{\mathrm{analytic}}\) is not an RH proof by itself]
\label{v4-arith:r-comparison:cor:hypothesis-also-wrong}
\ClaimStatusProvedHere. The equivalences in the chain
\(\mathrm{Det}\Leftrightarrow\mathrm{VB}^{*}\text{ main direction}
\Leftrightarrow\mathrm{RH}\) are:
\begin{enumerate}
\item Determinacy is equivalent to the \(\mathrm{VB}^{*}\)
\emph{converse}, not the main direction
(\Cref{v4-arith:r-comparison:prop:det-converse-not-main}).
\item Determinacy needs input beyond RH in this reduction
(\Cref{v4-arith:r-comparison:prop:det-strictly-finer}).
\end{enumerate}
Hence \(\mathrm{R}_{\mathrm{analytic}}\) is not, by itself, an RH proof
in the present implication chain.
\end{corollary}

\begin{proof}
\Cref{v4-arith:vb:thm:three-equivalences} and
\Cref{v4-arith:hpvbconv:thm:VBstar-converse-spectral-regularity}
are unambiguous on both the direction of the equivalence and the
need for determinacy input beyond RH in this reduction.
\end{proof}

\section{The Sharp Reduction of \(\mathrm{R}_{\mathrm{analytic}}\)}
\label{v4-arith:r-comparison:honest-reduction}

\begin{theorem}[Sharp reduction of \(\mathrm{R}_{\mathrm{analytic}}\)]
\label{v4-arith:r-comparison:thm:honest-reduction}
\ClaimStatusProvedHere. The term
\(\mathrm{R}_{\mathrm{analytic}}\) is the demand for an unconditional
two-sided moment asymptotic \(\mu_{2n}\asymp C^{n}n!^{2}\). It has the
following properties.
\begin{enumerate}
\item \emph{Form}: classical-analytic; \(\mathrm{R}_{\mathrm{analytic}}\)
is a two-sided moment asymptotic.
\item \emph{Position in the implication hierarchy}: under
P2, P3, \(\mathrm{H\text{-}P}^{\mathrm{Ar}}\),
\(\mathrm{R}_{\mathrm{analytic}}\Rightarrow\mathrm{Det}
\Leftrightarrow\mathrm{VB}^{*}\) converse; it is not an RH proof
directly.
\item \emph{Implication to RH}: the present reduction does not prove
\(\mathrm{R}_{\mathrm{analytic}}\Rightarrow\mathrm{RH}\).
\item \emph{Difficulty}: removing the unconditional
\((\log n)^{4n}\) gap requires density-of-zeros hypotheses for
\(\zeta\), an open analytic-number-theory problem of RH-adjacent
difficulty.
\item \emph{Unconditional Carleman construction via
Bombieri--Lagarias and Conrey}: inconclusive; the correct
Bertrand-series lower bound \(\sum 1/(n(\log n)^{2})\) converges
(\Cref{v4-arith:r-comparison:cor:inconclusive}).
\end{enumerate}
\(\mathrm{R}_{\mathrm{analytic}}\) is classical in form, not proved
here, insufficient by itself to prove RH in this reduction, and
RH-adjacent in difficulty.
The summary of Chapter~\ref{v4-arith:hpar-det}
(``classical-analytic, not RH-equivalent'') is accurate on the form
and implication axes and inaccurate on the difficulty axis.
\end{theorem}

\begin{proof}
(1) is \Cref{v4-arith:r-comparison:axis-distinction}.
(2) is \Cref{v4-arith:r-comparison:diag:chain}.
(3) is \Cref{v4-arith:hpar-det:rem:residual-vs-RH}.
(4) is \Cref{v4-arith:r-comparison:prop:closing-gap-is-density-RH}.
(5) is \Cref{v4-arith:r-comparison:cor:inconclusive}.
\end{proof}

\begin{remark}[precise statement]
\label{v4-arith:r-comparison:rem:what-delta-should-say}
\(\mathrm{R}_{\mathrm{analytic}}\) asks for an unconditional two-sided
moment asymptotic at the \(n!^{2}\) rate. Its form is
classical-analytic. Its implication relative to RH is one-directional
in this reduction. Its difficulty is RH-adjacent: removing the
unconditional Bombieri--Lagarias--Conrey gap requires
zero-density-hypothesis-level improvements. Carleman's criterion
via the Bombieri--Lagarias--Conrey bounds alone is inconclusive; a
genuine Carleman divergence requires either density-of-zeros
improvements or the unconditional two-sided asymptotic
\(\mu_{2n}\asymp C^{n}n!^{2}\) without logarithmic inflation. The
remaining work is zero-density analytic number theory.
\end{remark}

\begin{remark}[domain of the statement]
\label{v4-arith:r-comparison:rem:not-sharpen}
The three terms \(\mathrm{W}+\mathrm{H}+\mathrm{R}_{\mathrm{analytic}}\)
of \Cref{v4-arith:hpar-det:cor:wave3-structure} remain independent.
The Stirling computation affects only the Carleman-divergence inference.
The weight selection (W) and Weil formula at Gaussian test functions
(H) components are unchanged.
\end{remark}

\section{Consequences}
\label{v4-arith:r-comparison:summary}

\begin{enumerate}
\item The characterisation of \(\mathrm{R}_{\mathrm{analytic}}\) as
``classical-analytic, not RH-equivalent''
(\Cref{v4-arith:r-comparison:thm:conclusion}) is accurate on the form axis
and the implication axis.
The form is classical-analytic; the implication does not by itself
prove RH; but the difficulty is RH-adjacent via the density-of-zeros
problem.
\item The Bombieri--Lagarias and Conrey estimates give
\Cref{v4-arith:r-comparison:cor:inconclusive}. The lower bound
on the Carleman sum is the Bertrand series \(\sum 1/(n(\log n)^{2})\),
which converges. The Bombieri--Lagarias construction is inconclusive.
\item The form of the chain
\(\mathrm{Det}\Leftrightarrow\mathrm{VB}^{*}\text{ main direction}
\Leftrightarrow\mathrm{RH}\) is given by
\Cref{v4-arith:r-comparison:cor:hypothesis-also-wrong}: determinacy is
equivalent to the \(\mathrm{VB}^{*}\) converse
(\Cref{v4-arith:hpvbconv:thm:VBstar-converse-spectral-regularity}),
and determinacy needs input beyond RH in this reduction.
\item \Cref{v4-arith:r-comparison:thm:honest-reduction} separates three
axes: form (classical), implication direction (not an RH proof by
itself), and difficulty (RH-adjacent via density-of-zeros).
\(\mathrm{R}_{\mathrm{analytic}}\) is the single analytic input of
the sufficient path; proving it is a genuine open problem of analytic
number theory at the zero-density level.
\item The three terms
\(\mathrm{W}+\mathrm{H}+\mathrm{R}_{\mathrm{analytic}}\) of
\Cref{v4-arith:hpar-det:cor:wave3-structure} remain the sharp reduction
of Chapter~\ref{v4-arith:rh-reformulation}.
\item Proving \(\mathrm{R}_{\mathrm{analytic}}\) requires either
(a) a density-of-zeros improvement pushing the Conrey lower bound
from 40\% of zeros on the line to the density-1 hypothesis with
sub-logarithmic error, or (b) an alternative positivity or
measure-regularisation construction producing
\(\mu_{2n}\asymp C^{n}n!^{2}\) directly. Either construction is a
substantive analytic-number-theory theorem. The Vol~IV sufficient
argument for RH rests on three open problems of comparable difficulty
(W, H, \(\mathrm{R}_{\mathrm{analytic}}\)).
\end{enumerate}

The three inputs are genuine open problems of substantive
analytic-number-theory difficulty. The term
\(\mathrm{R}_{\mathrm{analytic}}\) is specifically a density-of-zeros
statement adjacent to RH.
